\documentclass[ug.tex]


\begin{document}

\chapter{Classical Dynamics}

\boxx{
\textbf{Classical Mechanics of Point Particles:}
\begin{itemize}
    \item Generalised Coordinates and Velocities.
    \item Hamilton's Principle, Lagrangian, and Euler-Lagrange Equations.
    \item Applications to Simple Systems Such as Coupled Oscillators.
    \item Canonical Momenta \& Hamiltonian.
    \item Hamilton's Equations of Motion.
    \item Applications: Hamiltonian for a Harmonic Oscillator, Particle in a Central Force Field.
    \item Poisson Brackets.
    \item Canonical Transformations. (25 Lectures)
\end{itemize}
}
\boxx{
\textbf{Special Theory of Relativity:}
\begin{itemize}
    \item Postulates of Special Theory of Relativity.
    \item Lorentz Transformations.
    \item Minkowski Space.
    \item The Invariant Interval, Light Cone, and World Lines.
    \item Space-time Diagrams.
    \item Time-dilation, Length Contraction \& Twin Paradox.
    \item Four-vectors: Space-like, Time-like \& Light-like.
    \item Four-velocity and Acceleration.
    \item Four-momentum and Energy-momentum Relation.
    \item The Electromagnetic Field Tensor and Its Transformation Under Lorentz Transformations.
    \item Relation to Known Transformation Properties of E and B.
    \item Electric and Magnetic Fields Due to a Uniformly Moving Charge.
    \item Equation of Motion of Charged Particle \& Maxwell's Equations in Tensor Form.
    \item Motion of Charged Particles in External Electric and Magnetic Fields. (35 Lectures)
\end{itemize}

\textbf{Electromagnetic Radiation:}
\begin{itemize}
    \item Review of Retarded Potentials.
    \item Potentials Due to a Moving Charge: Lienard-Wiechert Potentials.
    \item Electric \& Magnetic Fields Due to a Moving Charge: Power Radiated, Larmor’s Formula and Its Relativistic Generalisation. (15 Lectures)
\end{itemize}
}



\subsection*{I. Generalized Coordinates and Velocities}

\subsubsection*{Statements:}

\begin{enumerate}
    \item \textbf{Introduction to Generalized Coordinates:}
    \begin{itemize}
        \item Define generalized coordinates and explain their significance in describing the configuration of a system.
    \end{itemize}
    
    \item \textbf{Generalized Velocities:}
    \begin{itemize}
        \item Introduce generalized velocities as the time derivatives of generalized coordinates.
    \end{itemize}
\end{enumerate}

\subsection*{II. Hamilton's Principle, Lagrangian, and Euler-Lagrange Equations}

\subsubsection*{Statements:}

\begin{enumerate}
    \item \textbf{Hamilton's Principle:}
    \begin{itemize}
        \item Present Hamilton's principle as a fundamental principle of classical mechanics.
    \end{itemize}
    
    \item \textbf{Lagrangian Formulation:}
    \begin{itemize}
        \item Introduce the Lagrangian as a function that summarizes the dynamics of a system.
    \end{itemize}
    
    \item \textbf{Euler-Lagrange Equations:}
    \begin{itemize}
        \item Derive the Euler-Lagrange equations from Hamilton's principle and explain their significance in determining the equations of motion.
    \end{itemize}
\end{enumerate}

\subsection*{III. Applications to Simple Systems: Coupled Oscillators}

\subsubsection*{Statements:}

\begin{enumerate}
    \item \textbf{Application to Coupled Oscillators:}
    \begin{itemize}
        \item Apply the Lagrangian formulation and Euler-Lagrange equations to the analysis of coupled oscillators.
    \end{itemize}
\end{enumerate}

\subsection*{IV. Canonical Momenta and Hamiltonian}

\subsubsection*{Statements:}

\begin{enumerate}
    \item \textbf{Canonical Momenta:}
    \begin{itemize}
        \item Define canonical momenta and explain their role in the Hamiltonian formulation.
    \end{itemize}
    
    \item \textbf{Hamiltonian:}
    \begin{itemize}
        \item Introduce the Hamiltonian as the Legendre transformation of the Lagrangian and discuss its significance.
    \end{itemize}
\end{enumerate}

\subsection*{V. Hamilton's Equations of Motion}

\subsubsection*{Statements:}

\begin{enumerate}
    \item \textbf{Hamilton's Equations Overview:}
    \begin{itemize}
        \item Present Hamilton's equations of motion as an alternative formulation of classical mechanics.
    \end{itemize}
    
    \item \textbf{Applications: Harmonic Oscillator, Central Force Field:}
    \begin{itemize}
        \item Apply Hamilton's equations to derive the equations of motion for specific systems, such as a harmonic oscillator and a particle in a central force field.
    \end{itemize}
\end{enumerate}

\subsection*{VI. Poisson Brackets}

\subsubsection*{Statements:}

\begin{enumerate}
    \item \textbf{Introduction to Poisson Brackets:}
    \begin{itemize}
        \item Define Poisson brackets and explain their role in Hamiltonian mechanics.
    \end{itemize}
\end{enumerate}

\subsection*{VII. Canonical Transformations}

\subsubsection*{Statements:}

\begin{enumerate}
    \item \textbf{Overview of Canonical Transformations:}
    \begin{itemize}
        \item Introduce canonical transformations and their significance in preserving the form of Hamilton's equations.
    \end{itemize}
\end{enumerate}

\section*{Note to Teachers:}

\begin{itemize}
    \item Use illustrative examples and diagrams to explain the concepts of generalized coordinates, Lagrangian formulation, and Hamiltonian mechanics.
    \item Conduct practical exercises or demonstrations involving simple mechanical systems to enhance students' understanding.
    \item Encourage students to solve problems related to coupled oscillators, canonical momenta, and Hamiltonian mechanics to reinforce theoretical concepts.
    \item Relate the principles of classical mechanics to real-world applications, emphasizing the generality and versatility of the Hamiltonian approach.
\end{itemize}

\rule{\linewidth}{0.5mm}


subsection*{I. Postulates of Special Theory of Relativity}

\subsubsection*{Statements:}

\begin{enumerate}
    \item \textbf{Introduction to Special Theory of Relativity:}
    \begin{itemize}
        \item Define the key concepts of special relativity and its significance in understanding the behavior of objects moving at high speeds.
    \end{itemize}
    
    \item \textbf{Postulates:}
    \begin{itemize}
        \item Present and discuss the two postulates of special relativity, emphasizing the constancy of the speed of light and the principle of relativity.
    \end{itemize}
\end{enumerate}

\subsection*{II. Lorentz Transformations}

\subsubsection*{Statements:}

\begin{enumerate}
    \item \textbf{Derivation of Lorentz Transformations:}
    \begin{itemize}
        \item Derive the Lorentz transformations from the postulates, highlighting their role in relating coordinates between inertial reference frames.
    \end{itemize}
    
    \item \textbf{Minkowski Space:}
    \begin{itemize}
        \item Introduce Minkowski space as the mathematical framework for combining space and time.
    \end{itemize}
    
    \item \textbf{Invariant Interval, Light Cone, and World Lines:}
    \begin{itemize}
        \item Define the invariant interval, light cone, and world lines in the context of Minkowski space.
    \end{itemize}
    
    \item \textbf{Space-Time Diagrams:}
    \begin{itemize}
        \item Illustrate and explain the use of space-time diagrams in representing events and world lines.
    \end{itemize}
\end{enumerate}

\subsection*{III. Time-Dilation, Length Contraction, and Twin Paradox}

\subsubsection*{Statements:}

\begin{enumerate}
    \item \textbf{Time-Dilation and Length Contraction:}
    \begin{itemize}
        \item Explain the concepts of time-dilation and length contraction, providing examples and practical implications.
    \end{itemize}
    
    \item \textbf{Twin Paradox:}
    \begin{itemize}
        \item Present the twin paradox and discuss its resolution within the framework of special relativity.
    \end{itemize}
\end{enumerate}

\subsection*{IV. Four-Vectors and Four-Momentum}

\subsubsection*{Statements:}

\begin{enumerate}
    \item \textbf{Types of Four-Vectors:}
    \begin{itemize}
        \item Introduce space-like, time-like, and light-like four-vectors.
    \end{itemize}
    
    \item \textbf{Four-Velocity and Acceleration:}
    \begin{itemize}
        \item Define four-velocity and four-acceleration, emphasizing their relativistic nature.
    \end{itemize}
    
    \item \textbf{Four-Momentum and Energy-Momentum Relation:}
    \begin{itemize}
        \item Derive the relationship between four-momentum and energy-momentum, highlighting the conservation laws.
    \end{itemize}
\end{enumerate}

\subsection*{V. Electromagnetic Field Tensor and Maxwell's Equations in Tensor Form}

\subsubsection*{Statements:}

\begin{enumerate}
    \item \textbf{Introduction to Electromagnetic Field Tensor:}
    \begin{itemize}
        \item Define the electromagnetic field tensor and its components.
    \end{itemize}
    
    \item \textbf{Transformation Properties:}
    \begin{itemize}
        \item Discuss how the electromagnetic field tensor transforms under Lorentz transformations, relating it to the transformation properties of electric and magnetic fields.
    \end{itemize}
    
    \item \textbf{Electric and Magnetic Fields of a Moving Charge:}
    \begin{itemize}
        \item Derive the expressions for electric and magnetic fields produced by a uniformly moving charge.
    \end{itemize}
    
    \item \textbf{Equation of Motion of Charged Particle:}
    \begin{itemize}
        \item Present the equation of motion for a charged particle in electromagnetic fields.
    \end{itemize}
\end{enumerate}

\subsection*{VI. Motion of Charged Particles in External Electric and Magnetic Fields}

\subsubsection*{Statements:}

\begin{enumerate}
    \item \textbf{Lorentz Force Equation:}
    \begin{itemize}
        \item Introduce the Lorentz force equation and discuss how it governs the motion of charged particles in external fields.
    \end{itemize}
    
    \item \textbf{Maxwell's Equations in Tensor Form:}
    \begin{itemize}
        \item Present Maxwell's equations in tensor form, emphasizing their covariance under Lorentz transformations.
    \end{itemize}
\end{enumerate}

\section*{Note to Teachers:}

\begin{itemize}
    \item Use visual aids, diagrams, and animations to help students visualize concepts like Lorentz transformations and space-time diagrams.
    \item Encourage class discussions on the implications of time-dilation and length contraction, promoting critical thinking.
    \item Provide examples and exercises related to four-vectors, four-momentum, and the Lorentz force equation to reinforce theoretical concepts.
    \item Relate the principles of special relativity to modern technologies and experiments, showcasing its practical applications.
\end{itemize}

\rule{\linewidth}{0.5mm}


\subsection*{I. Review of Retarded Potentials}

\subsubsection*{Statements:}

\begin{enumerate}
    \item \textbf{Introduction to Retarded Potentials:}
    \begin{itemize}
        \item Provide a brief review of retarded potentials and their significance in understanding the electromagnetic interactions of charged particles.
    \end{itemize}
\end{enumerate}

\subsection*{II. Potentials due to a Moving Charge: Lienard-Wiechert Potentials}

\subsubsection*{Statements:}

\begin{enumerate}
    \item \textbf{Lienard-Wiechert Potentials Overview:}
    \begin{itemize}
        \item Introduce the Lienard-Wiechert potentials as solutions to the electromagnetic field equations for a moving charge.
    \end{itemize}
    
    \item \textbf{Expression for Lienard-Wiechert Potentials:}
    \begin{itemize}
        \item Derive the expressions for electric and magnetic potentials due to a moving charge, emphasizing the dependence on the velocity and distance.
    \end{itemize}
\end{enumerate}

\subsection*{III. Electric and Magnetic Fields due to a Moving Charge}

\subsubsection*{Statements:}

\begin{enumerate}
    \item \textbf{Electric Field from a Moving Charge:}
    \begin{itemize}
        \item Derive the electric field produced by a moving charge using the Lienard-Wiechert potentials.
    \end{itemize}
    
    \item \textbf{Magnetic Field from a Moving Charge:}
    \begin{itemize}
        \item Derive the magnetic field produced by a moving charge using the Lienard-Wiechert potentials.
    \end{itemize}
    
    \item \textbf{Power Radiated:}
    \begin{itemize}
        \item Discuss the concept of power radiated by a moving charge and its importance in understanding electromagnetic radiation.
    \end{itemize}
    
    \item \textbf{Larmor's Formula:}
    \begin{itemize}
        \item Introduce Larmor's formula as an expression for the power radiated by an accelerated charge.
    \end{itemize}
    
    \item \textbf{Relativistic Generalization of Larmor's Formula:}
    \begin{itemize}
        \item Extend Larmor's formula to its relativistic form, taking into account the effects of high velocities.
    \end{itemize}
\end{enumerate}

\section*{Note to Teachers:}

\begin{itemize}
    \item Utilize mathematical derivations and step-by-step explanations to make the concepts of retarded potentials and Lienard-Wiechert potentials accessible to students.
    \item Conduct demonstrations or simulations illustrating the behavior of electric and magnetic fields around a moving charge to enhance visualization.
    \item Engage students in problem-solving exercises related to Larmor's formula and its relativistic generalization to reinforce theoretical concepts.
    \item Discuss real-world applications and phenomena related to electromagnetic radiation, connecting theoretical knowledge to practical scenarios.
\end{itemize}



%------------Body-Ends--------------------
%\bibliography{ref.bib}
%\bibliographystyle{IEEEtran}
\end{document}
